% GENERATED FILE. Edit canonical lessons, then run npm run generate:books.
% canonical-source-hash: fnv1a64:55d8412b880c2ef3
% canonical-lessons: GE-C448-frankieren, GE-C448-verschicken, GE-C448-einwerfen, GE-C448-paeckchen, GE-C448-einschreiben

\chapter{Briefe und Päckchen}
\label{ch:ge-briefe-und-packchen}
\hlchaptermodality{448}

\begin{chapteropening}
\textbf{By the end of this chapter:} \emph{I can send things by post: putting a stamp on, sending, posting a letter, the small parcel and registered post.}

Complete the last lesson of chapter 448: I can send things by post: putting a stamp on, sending, posting a letter, the small parcel and registered post.
\end{chapteropening}

\section[frankieren]{frankieren — to put postage on, to stamp (a letter)}

\label{lesson:GE-C448-frankieren}



\begin{quote}
Before the new one: say the German for the watchmaker, then the German for the jeweller.
\end{quote}

\subsection*{You'll want to know: frankieren}
\textbf{frankieren} — \textquotedblleft{}to put postage on, to stamp (a letter)\textquotedblright{}.

\begin{quote}
Sie frankiert den Brief und bringt ihn zur Post.
\end{quote}

\begin{quote}
\emph{She puts a stamp on the letter and takes it to the post office.}
\end{quote}

Past participle without \emph{ge-}: \emph{ein frankierter Umschlag} — a stamped envelope.

\begin{cousinweb}
\emph{frankieren} was borrowed from Italian \emph{francare} 'to free from charges', built on \emph{franco} 'free': a stamped letter travels free. EtymDB does not list the German verb, so the line is hedged.
\end{cousinweb}

\begin{grammarlens}[title={What you've built}]
One more word for telling, explaining and describing at B1.
\end{grammarlens}

\subsection*{Your turn}
\begin{itemize}
\raggedright
  \item \emph{Say it:} \emph{frankieren}
  \item \emph{Say it:} \emph{frankieren}, once more
  \item \emph{Say it:} a sentence of your own with \emph{frankieren}
  \item \emph{Recall:} say the German for the watchmaker, then the German for the jeweller, then say \emph{frankieren} again
\end{itemize}

\begin{tcolorbox}[breakable,colback=teal!4,colframe=teal!35!black,title={Before you move on}]
What does frankieren mean? (\textquotedblleft{}To put postage on, to stamp (a letter)\textquotedblright{}.) Say it once more.
\end{tcolorbox}
\section[verschicken]{verschicken — to send (out)}

\label{lesson:GE-C448-verschicken}



\begin{quote}
Before the new one: say the German for the jeweller, then the German for to put postage on.
\end{quote}

\subsection*{You'll want to know: verschicken}
\textbf{verschicken} — \textquotedblleft{}to send (out)\textquotedblright{}.

\begin{quote}
Zu Weihnachten verschickt sie viele Karten.
\end{quote}

\begin{quote}
\emph{At Christmas she sends lots of cards.}
\end{quote}

\emph{Verschicken} often means sending many things, or sending something away.

\begin{cousinweb}
\emph{verschicken} is built from \emph{ver-} and \emph{schicken} 'to send', inherited from Middle High German \emph{schicken} 'to arrange, to put in order' (EtymDB).
\end{cousinweb}

\begin{grammarlens}[title={What you've built}]
One more word for telling, explaining and describing at B1.
\end{grammarlens}

\subsection*{Your turn}
\begin{itemize}
\raggedright
  \item \emph{Say it:} \emph{verschicken}
  \item \emph{Say it:} \emph{verschicken}, once more
  \item \emph{Say it:} a sentence of your own with \emph{verschicken}
  \item \emph{Recall:} say the German for the jeweller, then the German for to put postage on, then say \emph{verschicken} again
\end{itemize}

\begin{tcolorbox}[breakable,colback=teal!4,colframe=teal!35!black,title={Before you move on}]
What does verschicken mean? (\textquotedblleft{}To send (out)\textquotedblright{}.) Say it once more.
\end{tcolorbox}
\section[einwerfen]{einwerfen — to post (a letter); to throw in}

\label{lesson:GE-C448-einwerfen}



\begin{quote}
Before the new one: say the German for to put postage on, then the German for to send.
\end{quote}

\subsection*{You'll want to know: einwerfen}
\textbf{einwerfen} — \textquotedblleft{}to post (a letter); to throw in\textquotedblright{}.

\begin{quote}
Auf dem Weg zur Arbeit wirft er den Brief ein.
\end{quote}

\begin{quote}
\emph{On the way to work he posts the letter.}
\end{quote}

Separable, strong: \emph{wirft ein}, \emph{warf ein}, \emph{hat eingeworfen}. \emph{Der Briefkasten} — the postbox.

\begin{cousinweb}
\emph{einwerfen} is built from \emph{ein} 'in' and \emph{werfen} 'to throw', inherited from Old High German \emph{werfan} and Proto-Germanic \emph{werpaną} (EtymDB).
\end{cousinweb}

\begin{grammarlens}[title={What you've built}]
One more word for telling, explaining and describing at B1.
\end{grammarlens}

\subsection*{Your turn}
\begin{itemize}
\raggedright
  \item \emph{Say it:} \emph{einwerfen}
  \item \emph{Say it:} \emph{einwerfen}, once more
  \item \emph{Say it:} a sentence of your own with \emph{einwerfen}
  \item \emph{Recall:} say the German for to put postage on, then the German for to send, then say \emph{einwerfen} again
\end{itemize}

\begin{tcolorbox}[breakable,colback=teal!4,colframe=teal!35!black,title={Before you move on}]
What does einwerfen mean? (\textquotedblleft{}To post (a letter); to throw in\textquotedblright{}.) Say it once more.
\end{tcolorbox}
\section[das Päckchen]{das Päckchen — the small parcel, the packet}

\label{lesson:GE-C448-paeckchen}



\begin{quote}
Before the new one: say the German for to send, then the German for to post.
\end{quote}

\subsection*{You'll want to know: das Päckchen}
\textbf{das Päckchen} — \textquotedblleft{}the small parcel, the packet\textquotedblright{}.

\begin{quote}
Das Päckchen für meine Tante ist nicht schwer.
\end{quote}

\begin{quote}
\emph{The small parcel for my aunt is not heavy.}
\end{quote}

Same form in the plural. A small \emph{Paket}, cheaper to send.

\begin{cousinweb}
\emph{Päckchen} is built from \emph{Pack} 'pack, bundle', usually traced to Low German or Dutch \emph{pak}, and the ending \emph{-chen} 'little'.
\end{cousinweb}

\begin{grammarlens}[title={What you've built}]
One more word for telling, explaining and describing at B1.
\end{grammarlens}

\subsection*{Your turn}
\begin{itemize}
\raggedright
  \item \emph{Say it:} \emph{das Päckchen}
  \item \emph{Say it:} \emph{das Päckchen}, once more
  \item \emph{Say it:} a sentence of your own with \emph{das Päckchen}
  \item \emph{Recall:} say the German for to send, then the German for to post, then say \emph{das Päckchen} again
\end{itemize}

\begin{tcolorbox}[breakable,colback=teal!4,colframe=teal!35!black,title={Before you move on}]
What does das Päckchen mean? (\textquotedblleft{}The small parcel, the packet\textquotedblright{}.) Say it once more.
\end{tcolorbox}
\section[das Einschreiben]{das Einschreiben — the registered letter, registered post}

\label{lesson:GE-C448-einschreiben}



\begin{quote}
Before the new one: say the German for to post, then the German for the small parcel.
\end{quote}

\subsection*{You'll want to know: das Einschreiben}
\textbf{das Einschreiben} — \textquotedblleft{}the registered letter, registered post\textquotedblright{}.

\begin{quote}
Den Vertrag schickt er per Einschreiben.
\end{quote}

\begin{quote}
\emph{He sends the contract by registered post.}
\end{quote}

Same form in the plural. The verb \emph{sich einschreiben} means to enrol.

\begin{cousinweb}
\emph{Einschreiben} is the verb \emph{einschreiben} 'to write in, to register' used as a noun; \emph{schreiben} is an early borrowing from Latin \emph{scrībere} 'to write' (Lewis \& Short).
\end{cousinweb}

\begin{grammarlens}[title={What you've built}]
One more word for telling, explaining and describing at B1.
\end{grammarlens}

\subsection*{Your turn}
\begin{itemize}
\raggedright
  \item \emph{Say it:} \emph{das Einschreiben}
  \item \emph{Say it:} \emph{das Einschreiben}, once more
  \item \emph{Say it:} a sentence of your own with \emph{das Einschreiben}
  \item \emph{Recall:} say the German for to post, then the German for the small parcel, then say \emph{das Einschreiben} again
\end{itemize}

\begin{tcolorbox}[breakable,colback=teal!4,colframe=teal!35!black,title={Before you move on}]
What does das Einschreiben mean? (\textquotedblleft{}The registered letter, registered post\textquotedblright{}.) Say it once more.
\end{tcolorbox}
